% !TeX program = xelatex
\documentclass[11pt]{article}

\usepackage[a4paper, margin=1in]{geometry}
\usepackage{amsmath, amssymb}
\usepackage{listings}
\usepackage{xcolor}
\usepackage{hyperref}
\usepackage{microtype}
\usepackage{fontspec}
\usepackage{booktabs}
\usepackage{enumitem}

% Body font: Latin Modern Roman, matching the other notes (which load `lmodern`).
\setmainfont{Latin Modern Roman}
% Mono font must cover the Lean unicode used in code listings (⟨ ⟩ σ ≤ ∅ₜ ...).
% DejaVu Sans Mono has full coverage; Consolas lacks ⟨ ⟩ ∅ and would drop them.
\setmonofont{DejaVu Sans Mono}[Scale=0.85]

\sloppy
\setlength{\emergencystretch}{3em}
\setlist{itemsep=2pt, topsep=3pt}

% ---- shell / Lean code-listing style ------------------------------------
\lstdefinelanguage{shell}{
  morekeywords={python, lake, build, pip, install},
  sensitive=true,
  morecomment=[l]{\#},
  morestring=[b]"
}
\lstset{
  inputencoding=utf8,
  extendedchars=true,
  basicstyle=\ttfamily\footnotesize,
  keywordstyle=\color{blue!70!black}\bfseries,
  commentstyle=\color{gray}\itshape,
  stringstyle=\color{red!55!black},
  showstringspaces=false,
  breaklines=true,
  columns=fullflexible,
  frame=single,
  framesep=4pt,
  rulecolor=\color{black!30},
  backgroundcolor=\color{gray!6},
  xleftmargin=6pt,
  xrightmargin=6pt,
  aboveskip=0.7\baselineskip,
  belowskip=0.7\baselineskip
}
% Style for showing the tool's raw terminal output verbatim.
\lstdefinestyle{output}{
  language={},
  basicstyle=\ttfamily\scriptsize,
  keywordstyle={},
  commentstyle={},
  stringstyle={},
  backgroundcolor=\color{black!4},
}
% -----------------------------------------------------------------------

\newcommand{\code}[1]{\texttt{\detokenize{#1}}}

\title{\texttt{flagmatic\_to\_lean}: a user guide}
\author{}
\date{}

\begin{document}
\maketitle

\section{The tool and what you need}

\code{flagmatic_to_lean.py}\, converts a \textbf{Flagmatic certificate (JSON)} into a
\textbf{Lean proof file (\code{.lean})}. The certificate holds only \emph{numeric data}
--- matrices and a list of flags --- and this tool maps those numbers to Lean
identifiers and attaches the positive-semidefiniteness lemmas, density tables, and
product expansions automatically. A single final \code{lake build} then turns it into a
\textbf{theorem verified in the kernel}; the proof body is filled in automatically too.

\medskip\noindent\textbf{No data files to prepare.} Everything the proof needs is
generated \emph{inside Lean} when the file is built: the forbid-free flags, their
densities, and the product expansions are all emitted by \code{generate_*} commands the
tool writes into the \code{.lean}. There are \textbf{no \code{*.json} data files on disk},
no \code{FlagDef.lean} to edit, and no forbid graph to define by hand --- the tool writes
the forbidden graph out itself.

\medskip\noindent You need three things to start.

\begin{itemize}
  \item \textbf{Python 3 with NetworkX} --- runs the converter. Install NetworkX once with
        \code{pip install networkx}; the tool uses its graph atlas to enumerate the
        canonical graphs \emph{in memory} (this is how it resolves each flag to its Lean
        index without reading any JSON). This enumeration supports up to $n=7$ vertices.
  \item \textbf{A working Lean / \code{lake} setup and this repository} --- where the
        generated \code{.lean} is actually built and verified. The generic forbid
        machinery (\code{completeSym2Graph}, the \code{generate_*} elaboration commands, the
        proof tactics) all live in this repo.
  \item \textbf{One certificate JSON} --- the input to convert; it encodes \emph{which
        density inequality to prove under which forbid}. Ready-to-use examples are in
        \code{LeanFlagAlgebras/Flagmatic/Certificates/}.
\end{itemize}

\noindent\textbf{All commands below are run from the repository root.}

\section{The workflow, step by step}

The whole workflow is two converter commands followed by a build:

\begin{quote}
\textbf{\code{inspect} (sanity-check the certificate) $\;\to\;$ \code{gen-skeleton} (write
the \code{.lean}) $\;\to\;$ \code{lake build} (verify in the kernel).}
\end{quote}

\noindent The steps below use Mantel (forbid $K_3$, edge density $\le 1/2$) as the worked
example.

\subsection*{Step 1 --- \texttt{inspect}: check parsing and the forbid}

\begin{lstlisting}[language=shell]
python LeanFlagAlgebras/Flagmatic/flagmatic_to_lean.py inspect LeanFlagAlgebras/Flagmatic/Certificates/mantel_cert.json
\end{lstlisting}

\noindent If every \code{->} resolves without error, the certificate is good to go.
\code{inspect} resolves \emph{every} flagmatic string (raising on the first that fails), so
it doubles as a validity check. The \code{forbid graph:} line names the forbidden graph and
how it will be stated in Lean, and the final block prints the exact \code{generate_*}
commands \code{gen-skeleton} will emit --- there are no JSON files to locate.

\begin{lstlisting}[style=output]
=== mantel_cert.json ===
description: 2-graph; maximize 2:12 density; forbid 3:121323
bound: 1/2
forbid graph: 3:121323 = K_3  ->  def K3 : Sym2Graph 3 := completeSym2Graph 3   (Route A: complete)
              bound stated as  ≤[completeGraph (Fin 3)]  (ordinary forbid)

admissible graphs (host-size):
                  '3:'  ->  FlagAlgebra_3_0_0_0     density=0
                '3:12'  ->  FlagAlgebra_3_0_0_1     density='1/3'
              '3:1213'  ->  FlagAlgebra_3_0_0_2     density='2/3'

types:
        '1:'  ->  FlagType_1_0

sigma-flags per type:
  type 0 ('1:'):
                 '2:(1)'  ->  FlagAlgebra_2_1_0_0
               '2:12(1)'  ->  FlagAlgebra_2_1_0_1

generation commands (what gen-skeleton emits; no JSON on disk):
  def K3 : Sym2Graph 3 := completeSym2Graph 3
  generate_pruned_forbid_free_empty_typed_flags 2 K3
  generate_pruned_forbid_free_empty_typed_flags 3 K3
  generate_pruned_forbid_free_flags 2 1 0 K3
  generate_pruned_forbid_free_flags 3 1 0 K3
  generate_pruned_flag_pair_density_theorems 2 3 1 0 K3
  generate_pruned_forbid_free_mul_theorems 2 3 1 0 K3 (completeGraph (Fin 3)) (completeSym2Graph_finFlag_mem_forbiddenFlags 3)
\end{lstlisting}

\noindent\textbf{If you get stuck:} a \code{ValueError: not a flagmatic string} or a
\code{LookupError: no graph in graphs_N isomorphic ...} points at a malformed certificate
(a bad edge string, or a vertex count above the supported $n=7$). Fix the certificate and
re-run. Unlike the old pipeline, there is nothing to \emph{build first}: the flags are
generated at \code{lake build} time, so there is no missing-data state to repair here.

\subsection*{Step 2 --- \texttt{gen-skeleton}: generate the Lean proof file}

\begin{lstlisting}[language=shell]
python LeanFlagAlgebras/Flagmatic/flagmatic_to_lean.py gen-skeleton LeanFlagAlgebras/Flagmatic/Certificates/mantel_cert.json LeanFlagAlgebras/Flagmatic/Mantel.lean --namespace Mantel --force
\end{lstlisting}

\noindent This one line is the whole job. The file it writes contains the imports, the
\code{def K3} forbid graph and its \code{generate_*} commands, the matrices with their PSD
lemmas, the flag vectors, and the \textbf{automatically proved main theorem}.

\begin{lstlisting}[style=output]
wrote 5078 chars to ...\Mantel.lean (namespace Mantel)
\end{lstlisting}

\noindent Options: \code{--force} overwrites an existing file of the same name; omit
\code{--namespace} to derive it from the target file name; \code{--theorem-name} sets the
theorem name (default: derived from the certificate file name, e.g.\
\code{mantel_cert.json} $\to$ \code{mantel_flagAlgebra}).

\noindent For the supported certificate shapes the main theorem is filled in with no
\code{sorry}. (Only for a shape the converter does not yet handle --- e.g.\ a block whose
number of flags exceeds the built-in \code{Fin.sum_univ_*} table --- is the proof body left
as a \code{sorry} stub, with a comment saying so.)

\subsection*{Step 3 --- \texttt{lake build}: verify in the kernel}

\begin{lstlisting}[language=shell]
lake build LeanFlagAlgebras.Flagmatic.Mantel
\end{lstlisting}

\noindent If it finishes without errors, \textbf{the proof is verified in the Lean kernel}.
The result is \code{theorem mantel_flagAlgebra} inside the \code{Mantel.lean} you just
generated. This build is where the flags are actually enumerated, the densities computed,
and the matrices' positive-semidefiniteness re-checked --- so it is the step that does the
real work.

\noindent\textbf{If the build is very slow or runs out of memory:} that is almost always
the large \code{native_decide} / kernel-checking steps, not a timeout (the generated file
already sets \code{maxHeartbeats 0}). Build just this one module (as above), close other
heavy processes, and retry.

\medskip\noindent We used Mantel as the example throughout. \textbf{Any other certificate
works the same way} --- change only the certificate path and the target \code{.lean} file
name / namespace, and follow Steps 1--3 unchanged.

\section{Two kinds of forbid: complete and non-complete}

The tool picks its generation route automatically from the forbidden graph.

\paragraph{Complete-graph forbids ($K_3, K_4, K_5, \ldots$).} The common case. The forbid
is the library graph \code{completeSym2Graph r}, the tool emits \code{def K}$r$ and the
\code{generate_pruned_*} commands (as in the Mantel output above), and the theorem is stated
\code{≤[completeGraph (Fin r)]}. Nothing extra is needed for a new clique size --- $K_6$,
$K_7$, \ldots\ work the same way (subject to the $n \le 7$ enumeration limit).

\paragraph{Non-complete forbids ($C_4, C_5, P_k, \ldots$).} Also supported.
\code{gen-skeleton} writes the forbidden graph out explicitly from the certificate's edges
and uses the \emph{subgraph} generators:
\begin{lstlisting}[style=output]
def ForbidGraph : Sym2Graph 5 where
  edges := {s(0, 1), s(0, 4), s(1, 2), s(2, 3), s(3, 4)}
  edges_valid := by decide
generate_subgraph_free_empty_typed_flags 5 ForbidGraph
generate_subgraph_free_flags 5 1 0 ForbidGraph
generate_subgraph_free_flag_pair_density_theorems 3 5 1 0 ForbidGraph
generate_subgraph_free_mul_theorems 3 5 1 0 ForbidGraph
\end{lstlisting}
The theorem is stated \code{≤[ForbidGraph.toLabeledGraph.graph]}. This is fully automatic:
the \code{C5turan.lean} example in the repository is generated this way and builds with no
\code{sorry}. \code{inspect} reports the subgraph route directly:
\begin{lstlisting}[style=output]
forbid graph: 5:1223344551  ->  def ForbidGraph : Sym2Graph 5 where edges := {s(0, 1), s(0, 4), s(1, 2), s(2, 3), s(3, 4)}   (Route B: subgraph)
              bound stated as  ≤[ForbidGraph.toLabeledGraph.graph]  (subgraph forbid)
\end{lstlisting}

\section{Example certificates to try}

Practice Steps 1--3 with the certificates below, all in
\code{LeanFlagAlgebras/Flagmatic/Certificates/}. Each produces a complete, \code{sorry}-free
proof.

\begin{center}
\small
\begin{tabular}{lccccc}
\toprule
certificate & forbid & route & host size $N$ & blocks & bound \\
\midrule
\code{mantel\_cert}        & $K_3$ & complete   & 3 & 1 & $1/2$ \\
\code{K3forbidP3\_cert}    & $K_3$ & complete   & 3 & 1 & $3/4$ \\
\code{K3forbidC4\_cert}    & $K_3$ & complete   & 4 & 2 & $3/8$ \\
\code{ErdosPentagon\_cert} & $K_3$ & complete   & 5 & 3 & $24/625$ \\
\code{K3forbidC6\_cert}    & $K_3$ & complete   & 6 & 2 & $92129/5242880$ \\
\code{K4turan\_cert}       & $K_4$ & complete   & 4 & 2 & $2/3$ \\
\code{K5turan\_cert}       & $K_5$ & complete   & 5 & 4 & $3/4$ \\
\code{C5turan\_cert}       & $C_5$ & subgraph   & 5 & 4 & $1/2$ \\
\bottomrule
\end{tabular}
\end{center}

\noindent The first seven forbid a complete graph; \code{C5turan} forbids the $5$-cycle and
exercises the non-complete (subgraph) route.

\section{Reference: the other subcommands}

\code{gen-skeleton} is the one you normally use. Two append-only helpers exist for editing an
existing file rather than regenerating it:

\begin{itemize}
  \item \code{gen-matrices <cert>.json <target>.lean} --- append just the \code{M_t} /
        \code{dM_t} / \code{LM_t} definitions and the PSD theorem.
  \item \code{gen-vectors <cert>.json <target>.lean} --- append just the \code{σ_t} and
        \code{v_t} flag-vector definitions.
\end{itemize}

\noindent Per-command help is available with
\code{python flagmatic_to_lean.py <subcommand> --help}.

\end{document}
